\section{Prototype Implementation}\label{sec:impl}

\subsection{Three Tracks}\label{sec:env}
\textbf{Track A -- Solidity on an EVM.} \texttt{IdentityRegistry.sol}, \texttt{ZeroTrustAccess.sol} and \texttt{RevocationGovernance.sol} (Solidity 0.8.24, optimizer 200 runs) implement Section~\ref{sec:arch}: the status machine, Algorithm~1 with anomaly gate and strict mode, \texttt{revokeBatch}/\texttt{suspendBatch}, PEP-authorized root anchoring and quorum voting. They run on the Hardhat in-process EVM with a Mocha/Chai suite of 16 tests (Section~\ref{sec:assure}) that also covers the industrial additions: the gate denies a high composite score when $A>g$, a tampered anomaly value invalidates the oracle signature, strict resources never step up, batch revocation skips unknown entries and blocks later access, and only authorized PEPs can anchor. This track gives real gas but not multi-node consensus or network latency, and the device's request signature is verified by the PEP rather than on-chain.

\textbf{Track B -- Python testbed.} The same state machine uses real Ed25519, SHA-256 and a Merkle root, and contains the \texttt{FastPathPEP}, which verifies the device signature and the oracle attestation, applies the per-device counter and tier policy, and logs decisions in Merkle-rooted batches. It is used for the scale, baseline, tier and governance experiments on a 2-vCPU Linux container (Python 3.13); all seeds derive from 20261004. Track B uses Ed25519 and Track A secp256k1 ECDSA, so their timings are not comparable.

\textbf{Track C -- Four-validator QBFT network.} The contracts are deployed on Hyperledger Besu 25.7.0 with four validators (QBFT, 2\,s block period, 4\,s round timeout), first on one two-vCPU host and then with one e2-standard-2 virtual machine per validator in four Google Cloud regions (Section~\ref{sec:gcp_eval}). \texttt{besu\_benchmark.js} (ethers v6) measures block interval, call latency, revocation visibility on a second validator, throughput with pre-signed transactions from ten senders, quorum voting and batch recall, and \texttt{besu\_fault.js} stops and restarts validators. The scripts and provisioning files are in the repository.

\subsection{Industrial Fleet and Signal Model}\label{sec:devsim}
The synthetic fleet has six classes: sensor 40\%, PLC 15\%, HMI 8\%, robot 7\%, gateway 15\% and legacy 15\%. Native devices draw posture in $[0.75,1]$, anomaly in $[0,0.15]$ and network context in $[0.7,1]$; legacy devices draw posture in $[0.55,0.75]$, anomaly in $[0,0.20]$ and network in $[0.6,0.9]$. To mimic process upsets, 5\% of legitimate requests by default (0, 5 and 15\% in the sensitivity study) draw anomaly from $[0.15,0.45]$. Four attack populations are used: \emph{naive compromise} (all signals degraded), \emph{partial mimicry} (posture and network legitimate, $A\in[0.3,0.7]$), \emph{stealthy mimicry} ($A\in[0.15,0.40]$) and \emph{full spoof} ($P=N=1$, $A\in[0.3,0.7]$); each starts from three clean prior requests. The legacy gateway of Section~\ref{sec:legacy} is modelled through the class \texttt{legacy} and its weaker signals; deriving proxy DIDs from the gateway key is a design element we did not implement. These distributions are our assumptions, not field measurements, and the attack and legitimate ranges overlap by construction, so they show how the policy responds to such overlap rather than how well a detector would perform in a plant. The real-data experiment of Section~\ref{sec:real_eval} replaces only the behavioural signal.

\subsection{Reproducibility}\label{sec:scenarios}
Every experiment is a command in the repository (\texttt{python -m sim.benchmark}, \texttt{sim.experiments\_extra}, \texttt{sim.experiments\_governance}, \texttt{sim.experiments\_iiot}, \texttt{sim.experiments\_prior}, \texttt{sim.experiments\_realdata}, and \texttt{hardhat run scripts/measure\_gas.js} and its governance and industrial variants), with results written as JSON and \texttt{paper/gen\_tables.py} regenerating every table. The TON\_IoT and WUSTL files are not redistributed.
